% =====================================================================
%  DECLARACIÓN DE USO DE HERRAMIENTAS DE INTELIGENCIA ARTIFICIAL
%  OBLIGATORIA si has usado cualquier herramienta de IA generativa
%  (asistentes de chat, autocompletado de código, traducción, generación
%  de imágenes...). Si no has usado ninguna, dilo en una frase.
%  Sé concreto: herramienta, para qué, en qué partes y cómo lo revisaste.
%  Lo que NO se acepta como aportación propia: el análisis, la discusión
%  de resultados ni las conclusiones generadas por IA.
% =====================================================================
\chapter*{Declaración de uso de herramientas de IA}
\pdfbookmark[0]{Declaración de uso de IA}{declaracionia}

Durante la elaboración de este trabajo se han utilizado las herramientas de
inteligencia artificial generativa que se indican a continuación. Tras su uso,
el autor o la autora ha revisado y editado todo el contenido afectado y asume la
plena responsabilidad del contenido de esta memoria y del código entregado.

% Tabla sin número ni pie: no debe aparecer en el índice de tablas.
\begin{center}
  \small
  \begin{tabularx}{\textwidth}{@{}>{\raggedright\arraybackslash}p{2.6cm}LL>{\raggedright\arraybackslash}p{3.2cm}@{}}
    \toprule
    \textbf{Herramienta} & \textbf{Uso concreto} & \textbf{Partes afectadas} & \textbf{Revisión} \\
    \midrule
    Asistente X (versión) & Corrección ortográfica y de estilo & Capítulos 1--3 & Lectura completa de los cambios \\
    Asistente de código Y & Generación de pruebas unitarias & Módulo \texttt{reservas} & Ejecución y revisión de cada prueba \\
    Herramienta Z & Búsqueda inicial de fuentes & Capítulo 2 & Cada fuente leída y verificada en el original \\
    \bottomrule
  \end{tabularx}
\end{center}

% Si no se ha usado ninguna herramienta, sustituye todo lo anterior por:
% «En la elaboración de este trabajo no se ha utilizado ninguna herramienta de
% inteligencia artificial generativa.»

\cleardoublepage
